\documentclass[11pt]{article}
\usepackage[margin=1in]{geometry}
\usepackage{amsmath,amssymb,amsthm,booktabs}
\usepackage[hidelinks]{hyperref}
\newtheorem{proposition}{Proposition}
\newtheorem{lemma}{Lemma}
\newcommand{\R}{\mathbb{R}}
\title{When Is a Mixture of Experts Easier to Certify?\\
Route-Aware Reductions and a Capacity-Matched Study}
\author{Author details omitted for technical review}
\date{Working manuscript, 31 October 2026}
\begin{document}
\maketitle
\begin{abstract}
Conditional computation reduces the number of expert parameters used by a
single MoE execution. It need not reduce the cost of certifying every execution
on a perturbation domain. We formulate an architecture-level comparison that
separates total capacity, active capacity, complete routing coverage, and
verification cost. A route-aware reduction exposes all legal selected-softmax
top-2 histories to a general nonlinear verifier. A frozen study compares
ordinary-trained MoE with total- and active-parameter-matched dense models,
using three seeds, 4800 complete verification calls and 24000 incomplete-bound
evaluations. At $2/255$, the single-layer MoE has 114/600 numerical positives
against 115/600 and 116/600 for the two dense matches. Neither passes the
registered architecture screen; both difference intervals include zero.
CROWN-IBP yields 23/3000 versus 2/3000 positives, but with higher median cost.
Thus the results support a limited incomplete-bound tradeoff, not general
architectural dominance. The conditional certification-training study was not
started. Numerical backend positives and independently checked source proofs
remain different evidence levels.
\end{abstract}

\paragraph{Scope of this draft.}
The manuscript target is ICML 2027, but the official 2027 submission deadline
and page limit have not been verified. Eight pages is a working target, not a
claim of compliance. The completed verifier study is preserved in
\path{paper/review_main.tex}; its outcome tables and empirical limits are
incorporated in Appendix~\ref{app:empirical}, not a separate simultaneous submission. This draft does not
record human review, a clean empirical reproduction, or a submission.

\section{Question and hypotheses}
A sparse MoE can represent many expert functions while activating only a few
at an input. Certification asks a different question: can every legal execution
on a neighborhood be shown to preserve the output property? The routers,
alternative histories, and their construction costs cannot be ignored.
We distinguish three claims and their required evidence.
\begin{itemize}
\item C1: guarded static verification obligations cover the route-variable
network. The reduction below establishes this statement under explicit source
and bound-validity assumptions. Small routing overhead is an empirical
condition, not part of the theorem.
\item C2: at matched capacity, MoE obtains at least as much certification with
lower complete-request cost. Total and active parameter matching answer
different questions; both must be reported.
\item C3: a routing-aware training objective improves the joint tradeoff among
clean accuracy, certified accuracy, and verification cost relative to dense
certification training. This is a gated research hypothesis, not a consequence
of a routing-margin penalty alone.
\end{itemize}
The contribution is not a claim that hybrid zonotopes, case splitting, interval
training, weak duality, or endpoint minimization are individually new.
Precedents already study sparse-MoE robustness~\cite{robmoe}, compositional
MoE verification~\cite{metamoe}, and certification with routed smoothed
experts~\cite{dualrs}. The new empirical question needs a shared deterministic
function, threat model, capacity convention, and full cost denominator.

\section{Route-aware complete-output verification}
At layer $\ell$, let $S_\ell=\{a,b\}$ be an unordered selected pair and let
$h_{\ell-1}$ be the output of its preceding execution. Define
\[
 m_\ell=r_{\ell,a}(h_{\ell-1})-r_{\ell,b}(h_{\ell-1}),\qquad
 h_\ell=\sigma(m_\ell)E_{\ell,a}(h_{\ell-1})+
        (1-\sigma(m_\ell))E_{\ell,b}(h_{\ell-1}).
\]
All selected scores must be at least all unselected scores. Ties admit every
legal pair. For a history $H=(S_1,\ldots,S_L)$, $F_H$ is the specialized graph
and $g_{H,j}(x)\geq0$ its prefix-dependent legality conditions. On input set
$X$, write $X_H=\{x\in X:g_{H,j}(x)\geq0\text{ for all }j\}$ and
$p_{H,k}(x)=q_k^\top F_H(x)+c_k$. We seek
\[
 \forall H\;\forall x\in X_H\;\forall k,\qquad p_{H,k}(x)>0.
\]

\begin{proposition}[Complete history reduction]
If a finite history cover contains every legal execution on $X$, and each
history is proved infeasible or every output property has a positive valid
lower bound over $X_H$, the complete MoE satisfies all output properties on $X$.
\end{proposition}
\begin{proof}
Any legal execution has a covered history, satisfies its guards, and agrees
with its specialized graph. Apply the lower bounds to that input. The argument
applies to every tie resolution, not just a library's selected clean route.
\end{proof}

\begin{lemma}[Prefix exclusion]
Suppose a proposed next pair requires $r_a-r_c\geq0$. A valid strict negative
upper bound on $r_a-r_c$, over a domain containing every legal execution of the
preceding prefix, excludes this extension and all its descendants.
\end{lemma}
\begin{proof}
The required inequality is impossible on that prefix. Any descendant requires
the same extension. A zero upper bound does not exclude equality.
\end{proof}
Our frontend computes its own prefix bounds; it does not borrow another arm's
route census. Numerical exclusions are recorded as such. A surviving candidate
is not a reachability witness, and an unfinished tree cannot establish safety.

\begin{proposition}[Guarded properties as output atoms]
For a fixed history and property,
\[
 \forall x\in X_H:p_H(x)>0\quad\Longleftrightarrow\quad
 \forall x\in X:\max\{p_H(x),-g_{H,1}(x),\ldots,-g_{H,J}(x)\}>0.
\]
\end{proposition}
\begin{proof}
At legal points every negated guard is nonpositive, so the maximum is positive
exactly when the property is positive. At illegal points a negated guard is
strictly positive. This also explains why a non-strict acceptance test would
incorrectly dismiss a violated property at a tie.
\end{proof}
ReLU implements the maximum, permitting a standard nonlinear verifier such as
GenBaB~\cite{genbab}. Logical equivalence does not imply that an output-side
guard produces the same relaxation as constraining the input domain.

\begin{proposition}[Last-gate endpoints]
If the suffix is affine and the upstream domain $P$ and retained guards do not
depend on the last independent gate $\lambda\in[l,u]$, then
\[
 \inf_{z\in P,\,\lambda\in[l,u]}[b(z)+\lambda d(z)]
 =\min\{\inf_{z\in P}[b(z)+l d(z)],\inf_{z\in P}[b(z)+u d(z)]\}.
\]
\end{proposition}
\begin{proof}
An affine function of $\lambda$ attains its minimum at an endpoint; infimum
commutes with the minimum of these two functions.
\end{proof}
Replacing gates by independently enclosing ranges is an outer relaxation,
not an equivalent network transformation. Every actual execution embeds by
using its actual gate values. Endpoints do not remove variable inner-gate
products or apply through a nonlinear suffix. This optional construction is
part of the prior study; the frozen S1 complete baseline uses coupled GenBaB.

\begin{lemma}[Why active parameters do not determine total cost]
For a fixed history only its selected expert graphs need be instantiated in
the specialized output, but its guards still require relevant router scores.
A complete request can incur
\[
 T=T_{\rm load}+T_{\rm route}+\sum_{H\in\mathcal H_{\rm visited}}
 (T_{\rm graph,H}+T_{\rm bound,H})+T_{\rm evidence/cleanup}.
\]
Thus bounding active parameters per history alone gives no bound independent
of the visited-history count or proof difficulty.
\end{lemma}
\begin{proof}
Static specialization removes unselected expert evaluations from that history;
it does not discharge other histories. Summing the work of the actually visited
obligations, including shared and non-shared overhead as implemented, gives the
decomposition. Arbitrarily many unresolved histories or expensive selected
networks remain consistent with the same per-history active parameter count.
\end{proof}
The sum describes work in one request. Parallel throughput across requests
must not be substituted for per-request cost. Our loader reads complete model
weights; that cost is charged even for weights absent from a specialized graph.

\section{Experiments}
\subsection{S1: frozen architecture screen}
CIFAR-10 uses a shared seed-0 45k/5k training/validation split, ToTensor only,
50 epochs, AdamW at 0.001 with weight decay 0.0001 and batch 256. Three seeds are
retained for each model. Earliest maximal validation accuracy selects the
checkpoint; test accuracy does not select it. Existing L1E8/L2E8 models use
cross-entropy plus their original 0.1 routing-balance term; dense controls use
cross-entropy. This architecture-specific objective difference is disclosed,
not presented as a pure architecture intervention with identical losses.

L1E8 has a 3072--128--8 router and eight 3072--256--128--10 experts. Router
parameters count in both capacity conventions. The dense width ratio 2:1 and
nearest integer parameter count are fixed before outcomes, without dummy
weights.
\begin{center}
\begin{tabular}{lrrr}
\toprule
Matching convention&MoE target&Dense parameters&Dense hidden widths\\
\midrule
Total&6,961,368&6,957,295&1758,879\\
Active&2,036,124&2,034,469&602,301\\
\bottomrule
\end{tabular}
\end{center}
Both mismatches are below 0.1\%. L2E8 is a secondary depth extension, not a
capacity match to these L1 controls. Logical active capacity is distinct from
training FLOPs: the old training code evaluates all experts before selection.

Evaluation uses the first 1000 test inputs for incomplete CROWN-IBP and first 200
for complete verification, at 2/255 and 8/255. Clean-wrong inputs remain in every
denominator. These deterministic prefixes have historical exposure; they are
development data, not a fresh independent holdout. Each complete request gets
300 seconds for startup, loading, all routing/graph/bound work, replay,
serialization, reception, and cleanup. The six MoE and six dense models give
4800 complete calls and 24000 incomplete items. At most eight single-thread CPU
workers run concurrently; each has an 8 GiB address-space limit. Both families
use the pinned numerical engine and common physical input boxes/properties;
dense uses native ReLU branching and MoE the coupled nonlinear configuration.

Public PGD uses 20 steps, two random starts and step size $\epsilon/4$, with
original-model replay. Its separately metered outcomes are not free answers
for either verifier. Numerical positives, replayed violations, unresolved
queries, timeouts and errors remain separate. No-witness attacks and negative
relaxation bounds do not prove safety and unsafety, respectively.

The primary endpoint is the 2/255 complete positive rate and all-request median
time for each L1 match. The frozen screen proceeds if either match has MoE rate
at least dense and a lower median time. Input-block bootstrap uses 10000
resamples with all three seed pairs retained. Intervals are descriptive here,
not an additional gate. A literal zero/zero rate comparison can pass this
screen on time; such a pass authorizes S2 but does not establish C2. All states,
missing costs and shared-clean-correct explanatory strata are reported.
Neither L2 nor 8/255 rescues a failed primary comparison.

\subsection{S1: completed training and attack audit}
All six dense controls completed their 50 frozen epochs. Validation-selected
test accuracies for seeds 0, 1, and 2 are respectively
$52.66\%,53.29\%,53.73\%$ for total matching and
$52.78\%,53.57\%,53.47\%$ for active matching. These are full-test clean
accuracies, not certification rates; all checkpoints remain in the comparison.
Training took 7998.453 seconds with one CPU worker. We do not compare that
training time to the earlier MoE runs, which used a different thread count.

The public attack queue settled all 2400 ten-input batches. Of these, 2392
were accepted; eight original storage-limit-period errors remain errors.
The complete domain denominator is 24000, including 80 domains without an
accepted batch result. Their partial saved records are not promoted offline.
The storage-only continuation increased the aggregate artifact cap from
128 to 512 MiB and retained all previous charges; neither the attack nor the
verifier configuration, input set, or request deadline changed.

\begin{center}
\begin{tabular}{lrrrr}
\toprule
Group&Clean correct&CE at $2/255$&Missing at $2/255$&CE at $8/255$\\
\midrule
L1E8&1492&2161&20&2921\\
L2E8&1386&2272&20&2953\\
Dense total&1632&2099&20&2944\\
Dense active&1637&2174&20&2972\\
\bottomrule
\end{tabular}
\end{center}
Every cell has denominator 3000: 1000 shared inputs and three fixed seeds,
not 3000 independent images. CE denotes an online full-model-replayed
counterexample record, including a clean misclassification. Clean counts use
the complete $8/255$ records; there are no missing domains or attack timeouts
at that radius. At $2/255$, the remaining $3000-\mathrm{CE}-20$ cases completed
the fixed search without a witness. The record audit checks bindings and online
acceptance, rather than rerunning every attack. Its summed request cost is
21107.093 seconds, not parallel wall time.

At $8/255$, only $0.93\%$--$2.63\%$ of these ordinarily trained model--input
pairs lack a replayed counterexample. This describes the difficulty of the
larger-radius task, not certified accuracy or an architecture advantage.
No-witness results remain unproved. All verification requests have now settled.
The final frozen-source record audit passed in 51.809 seconds (54.298 seconds
including process setup and publication), without new model or solver queries.
It retains all eight earlier PGD errors and reports zero new errors, unexecuted
verification requests, or missing request costs. Sources and audit identities
are recorded in \path{docs/architecture_s1_result_20261031.json}.

\subsection{S1 complete results: the registered screen stops}
Table~\ref{tab:s1primary} reports the primary endpoint. Neither capacity match
passes: total matching has one fewer MoE positive despite a slightly shorter
median, and active matching has two fewer with a slightly longer median.
The corresponding rate intervals include zero. This is a failed prespecified
screen, not a finding that MoE is universally inferior. The denominator is
200 shared inputs times three fixed seeds, not 600 independent images.

\begin{table}[ht]
\centering\small
\begin{tabular}{lrrrr}
\toprule
Dense match & Positives M/D & Gain/loss & $\Delta$ pp [95\% interval] & Median s M/D\\
\midrule
Total & 114/115 & 55/56 & $-0.167$ [$-3.833$,3.500] & 10.847/10.945\\
Active & 114/116 & 52/54 & $-0.333$ [$-3.833$,3.167] & 10.847/10.716\\
\bottomrule
\end{tabular}
\caption{Complete numerical positives at $2/255$, each arm out of 600.
M/D denotes MoE/dense; intersections are 59 and 62. All-request costs and
descriptive input-block intervals retain every seed and outcome.}
\label{tab:s1primary}
\end{table}

Seed-wise positive counts are $(46,34,34)$ for L1, $(26,47,42)$ for total
matching, and $(40,35,41)$ for active matching. The unmatched L2 extension has
$(4,7,6)$, or 17/600. Table~\ref{tab:s1states} retains all complete-verification
states. A timeout or unknown is never relabeled unsafe.

\begin{table}[ht]
\centering\small
\begin{tabular}{llrrrrr}
\toprule
Radius & Group & Positive & Replayed & Unknown & Timeout & Mean s\\
\midrule
$2/255$ & L1 &114&312&28&146&92.759\\
 & Dense total &115&275&133&77&105.015\\
 & Dense active &116&272&191&21&87.520\\
 & L2 &17&346&3&234&121.939\\
$8/255$ & L1 &0&312&3&285&144.455\\
 & Dense total &0&275&177&148&161.590\\
 & Dense active &0&272&282&46&162.896\\
 & L2 &0&346&0&254&128.579\\
\bottomrule
\end{tabular}
\caption{All complete requests, 600 per row; no verification errors or missing
calls. Replayed includes clean misclassification. The eight retained errors
belong to auxiliary PGD. Means include every request.}
\label{tab:s1states}
\end{table}

Cost medians need care. At $2/255$, the shared-clean-correct explanatory strata
contain 207 pairs for total matching, with medians 289.060/242.424s (MoE/dense),
and 217 for active matching, with 289.451/84.162s. These are separately recorded
strata, not replacements for the primary all-input denominator. The all-request
median-difference intervals are [$-1.576$,0.428]s and [$-2.539$,0.499]s.

At $8/255$, all positive counts are zero, including incomplete verification.
The complete all-request median is 11.735s for L1 and 289.485/289.375s for dense.
MoE has 312 clean-wrong replays, over half the denominator, while dense has
275/272, below half. The median therefore switches from cheap rejection to
difficult verification. On shared-clean-correct pairs all three medians are
approximately 290s. This is not a 25-fold useful certification speedup, and the
secondary radius cannot rescue the failed primary gate.

\subsection{Incomplete bounds: a coverage--cost tradeoff}
At $2/255$, incomplete CROWN-IBP~\cite{crownibp} produces 23/3000 L1 positives
against 2/3000 for each dense match: 22 gains, one loss, and one common positive.
The rate difference is 0.700 percentage points, descriptive 95\% interval
[0.400,1.067]. L1 median cost is 7.588s, versus 7.251s (total) and 7.140s (active):
4.64\% and 6.27\% higher. L2 has no incomplete positive. This is evidence of a
limited incomplete-bound advantage with a cost penalty, not of the primary
complete-verification hypothesis. The 1000 input blocks retain all three seeds;
these results neither replace the primary endpoint nor authorize S2.

Complete verification costs 167.459 summed request hours, incomplete 51.329,
and public PGD 5.863: 224.651 hours total. The eight-worker parallel queue wall
time is 28.094 hours across two segments, excluding the approval gap, training
and audits. Training costs 2.222 hours separately. Parallel throughput is not
per-request speed. Full states, paired gain/loss identities, cost intervals
and source hashes are archived. Route counts are conservative numerical
candidate upper counts, not exact legal-history counts or route-flip witnesses.

\subsection{S2 and ablations: not executed after the stop}
The registered rule returns this study to reductions and empirical limits.
Certification training, routing regularization and their proposed ablations
were not executed. In particular, evaluating ordinary-trained networks using
CROWN-IBP does not test CROWN-IBP training or the C3 tradeoff claim. The failed
screen does not prove such training could never help, but it does not justify
starting it under this protocol. No extra samples, new favorable radius, or
specialized-verifier mechanism was added to change the decision.

\subsection{Earlier verifier experiments as motivation}
The archived multilayer study obtains 108/180 numerical positives with pruned
GenBaB and none with ACT v1. A subsequent gate-decoupling cascade obtains 122/300
against 116/300: nine gains, three losses, net 2 percentage points, input-block
interval [0,4]. It fails its five-point threshold. The corrected guard diagnostic
closes one of eleven histories and fails admission. These are retained negative
results, not evidence that dense architectures are harder to certify or that
expert ReLUs uniquely explain every unresolved query. The cohorts were selected
differently from S1 and cannot be pooled into its accuracy denominator.
Appendix~\ref{app:empirical} incorporates the archived model and outcome tables,
their selection boundaries, and the stopped mechanisms; these are not new runs.

\section{Related work and limits}
Sparse-MoE robustness and dense compute comparisons precede this study~\cite{robmoe}.
MetaMoE uses class-separated hard top-1 compositional verification~\cite{metamoe}.
Dual randomized smoothing certifies a different, statistical $L_2$ function
with routed experts~\cite{dualrs}. RT-ER/JTDMoE already motivates expert-focused
robust training~\cite{rter}. CertMoE is recent route-aware text certification
work~\cite{certmoe}. None of these differences warrants a blanket ``first MoE
certification'' claim, and attack accuracies or smoothed-function guarantees
cannot substitute for our deterministic $L_\infty$ endpoint. The five-work source
review is available in the project record; author claims are not independent
proofs by this study.

This study covers three fixed seeds of a small ordinary-trained MLP family,
approximately matched parameter counts, historically exposed input prefixes,
and an architecture-specific balance loss. It is not a controlled theorem
about all sparse architectures or a test of certification-trained networks.
All experimental backend positives use a floating numerical policy. Exact
arithmetic on a supplied LP cannot repair an unverified input enclosure or
network conversion. Historical source gaps remain, and the new study does
not claim a source-complete proof of native PyTorch execution. The mathematical
reductions assume valid sources and bounds. A successful record audit checks
identity and coverage, not every numerical bound. Human technical review and
clean-environment reproduction remain outstanding.

\section{Disposition}
S1 is complete and fails the registered architecture screen under both
capacity conventions. C2 is not established, and C3 was not tested. The
paper's supported content is route-aware reductions, a bounded empirical
comparison, and the limits of active-parameter and aggregate-timing arguments.
The incomplete-bound gain remains part of the result. It is not used to
override the stop or to revive an unsuccessful mechanism. Independent human
review and clean empirical reproduction are still required.

\begin{thebibliography}{99}
\bibitem{robmoe} J. Puigcerver et al. On the Adversarial Robustness of Mixture
of Experts. NeurIPS 2022. \url{https://doi.org/10.52202/068431-0702}.
\bibitem{metamoe} Q. Pham et al. MetaMoE: Formal Verification of Compositional
Robustness and Scalability of Mixture-of-Experts Architecture. SAIV 2026.
\url{https://doi.org/10.1007/978-3-032-32357-6_8}.
\bibitem{dualrs} C. Sun, Y. Mao, and M. Vechev. Dual Randomized Smoothing:
Beyond Global Noise Variance. ICLR 2026.
\url{https://arxiv.org/abs/2512.01782v4}.
\bibitem{genbab} Z. Shi et al. Neural Network Verification with Branch-and-Bound
for General Nonlinearities. TACAS 2025. \url{https://arxiv.org/abs/2405.21063}.
\bibitem{ibp} S. Gowal et al. On the Effectiveness of Interval Bound Propagation
for Training Verifiably Robust Models. \url{https://arxiv.org/abs/1810.12715}.
\bibitem{crownibp} H. Zhang et al. Towards Stable and Efficient Training of
Verifiably Robust Neural Networks. ICLR 2020. \url{https://arxiv.org/abs/1906.06316}.
\bibitem{rter} X. Zhang, K. Xu, Z. Hu, and R. Wang. Optimizing Robustness and
Accuracy in Mixture of Experts: A Dual-Model Approach. ICML 2025.
\url{https://proceedings.mlr.press/v267/zhang25cj.html}.
\bibitem{certmoe} Q. Wang et al. CertMoE: Robustness Certification for
Mixture-of-Experts LLMs via Routing-Aware Smoothing. arXiv preprint 2610.29395, 2026.
\url{https://arxiv.org/abs/2610.29395}.
\end{thebibliography}
\clearpage
\appendix
\section{Archived verification study and limits}
\label{app:empirical}
This appendix incorporates the completed verifier-study evidence into the
current manuscript. It does not restart that study, pool its cohorts with S1,
or turn its numerical positives into source-complete proofs. The seven tables
below are included unchanged from the committed, reproducible summary view
used by the earlier technical-review draft.

\paragraph{Models, cohorts, and costs.}
In configuration names, $L$ is routed depth, $E$ the number of experts per
layer, and $R$ denotes a residual model. Routers have a width-128 ReLU hidden
layer; expert hidden widths are 256 and 128. Sequential networks flatten the
image and use width-128 intermediate routed outputs. Residual models instead
use a linear image-to-128 stem, blocks $h+\operatorname{MoE}(h)$, and a linear
ten-logit head. Both omit batch normalization. Seeds, split, epochs, optimizer,
and checkpoint selection are frozen. Sequential training averages the balance
penalty across layers, while residual training sums it; this is not a causal
ablation of residual connections alone.

The cascade study selects twenty common clean-correct inputs across all
eighteen models from 435 nonexcluded candidates. Its primary comparison uses
the fifteen multilayer models, hence 300 paired model--input objects on only
twenty input blocks. Its conditional positive rate is not S1's whole-prefix
certified accuracy. S1 uses the separate first-200/1000 all-input denominators,
including clean errors and disclosed historical exposure. Neither model nor
input selection may be transferred between these cohorts after outcomes.

Within the historical paired comparison both arms use the same disclosed R7
compatibility snapshot, CPU/double precision, one numerical thread, 8GiB
address-space limit, and 300-second full-request budget. GenBaB uses optimized
alpha bounds and nonlinear branching, with separately charged prefix bounds.
The results describe this configured engine, not every possible configuration.
All 1080 registered B4 calls were audited, with 56.319 charged execution hours;
the 240 unadmitted M-B/M-C ablation slots remain unused. Errors and timeout
costs stay in the denominator. Shared-host variation limits fine timing claims.

\paragraph{What the stopped tests establish.}
The earlier multilayer study has 600 calls: 540 across three configurations,
three seeds, twenty inputs and three arms, plus sixty single-layer pruning
anchors. Pruned GenBaB obtains 108/180 multilayer positives. ACT v1's zero
consists of 178 timeouts, one error and one unsupported result, not 180 solver
timeouts. A supplemental attribution gate registered after a 189-call prefix
finds many-history dominance in 2/52 L3E8 and 0/58 L2E16 objects after removing
replayed unsafe cases. Its timing is disclosed; this is not a fully
preregistered rejection of every possible routing-complexity bottleneck.

M-A replaces gates by independent enclosing intervals, checks last-gate
endpoints, and falls back to coupled solving with the remaining request time
after a phase cap of 120 seconds. Its nine gains all use the relaxed path;
48/52 endpoints close in initial propagation. Three complete-request losses
remain, so net gain is two percentage points with input-block interval [0,4],
below the registered five-point threshold. Gains concentrate in two-layer
models. The paid repeat of the coupled arm changes four of 120 terminal states,
including one positive-to-timeout change: that is not a three-percentage-point
estimate of positive-result noise. Neither a shared positive nor a retained
candidate refunds time spent on an unsuccessful cascade phase.

M-B uses nonnegative guard multipliers and weak duality. The original
independent-trajectory comparison could not isolate incremental guard benefit.
The revised test retains an evaluated twenty-step alpha prefix, then compares
forty guard/alpha updates against paid extra-alpha updates on the same eleven
histories. It adds one complete local closure, below the required five.
Original M-C partitions the first-layer margin of twenty-seven local histories
and adds no closures. A proposed last-layer dual-conditioned variant was gated
on revised M-B and was not executed. These outcomes do not prove that every
guard, margin partition or expert-network representation is ineffective.

The larger-radius and single-layer results below do not overturn those stops.
No gain appears among the thirteen larger-radius objects with multiple observed
histories; a single observed history elsewhere does not prove route stability.
The single-layer cohort has the same 57 positives in both arms, with extra
cascade cost. None of these diagnostic observations isolates ReLU difficulty
as the unique cause of every unresolved request.

% Relevant committed summary primary sha256 30c1121e079619babae3fb0bd0d6f09652d1e1be3ca0efb126b35342b94e875d
% Relevant committed summary multilayer sha256 5e78f74054d0dd883fbfe692c8a5146f155786fb917f0c491c54be8a49c6d127
% Relevant committed summary guard sha256 8bebc40ef4f01d9630a61de608044f2d730faf7b2574acf165b5262afa56f8fa
\begin{table}[ht]
\centering\small
\caption{Complete-request cascade: 300 calls per arm. P is numerical POSITIVE, not a strict certificate. All terminals and costs are retained.}
\label{tab:primary}
\begin{tabular}{lrrrrrr}
\toprule
Method & P & UNKNOWN & TIMEOUT & ERROR & Mean s & Median s \\
\midrule
Coupled & 116 & 8 & 176 & 0 & 189.287 & 288.350 \\
Cascade & 122 & 5 & 171 & 2 & 192.983 & 288.337 \\
\bottomrule
\end{tabular}
\end{table}

\begin{table}[ht]
\centering\small
\caption{Family breakdown of the same primary cohort: 60 model--input objects per row. Gain/loss counts are paired, not differences between unmatched totals.}
\label{tab:families}
\begin{tabular}{lrrrr}
\toprule
Family & Coupled P & Cascade P & Gained & Lost \\
\midrule
L2E8 & 24 & 28 & 5 & 1 \\
L3E8 & 27 & 26 & 0 & 1 \\
L2E16 & 53 & 53 & 0 & 0 \\
R2E8 & 7 & 10 & 4 & 1 \\
R3E8 & 5 & 5 & 0 & 0 \\
\bottomrule
\end{tabular}
\end{table}

\begin{table}[ht]
\centering\small
\caption{Separate multilayer bottleneck study: 60 objects per configuration/arm. The single-layer cost anchor is a separate pruning-only 60 calls.}
\label{tab:multilayer}
\begin{tabular}{lrrr}
\toprule
Configuration & ACT v1 P & Naive P & Pruned P \\
\midrule
L2E8 & 0 & 0 & 28 \\
L3E8 & 0 & 0 & 33 \\
L2E16 & 0 & 0 & 47 \\
\bottomrule
\end{tabular}
\end{table}

\begin{table}[ht]
\centering\small
\caption{Revised guard diagnostic on the same eleven histories. Each continuation retains its complete twenty-evaluation baseline. A local history is not a full MoE request.}
\label{tab:guard}
\begin{tabular}{lrr}
\toprule
Continuation & Closed histories & UNKNOWN \\
\midrule
Extra alpha & 0 & 11 \\
Joint alpha/guard & 1 & 10 \\
\bottomrule
\end{tabular}
\end{table}

\clearpage

\section{S1 seed-level counts and evidence reconstruction}
\label{app:s1seeds}
Table~\ref{tab:s1seeds} preserves all twelve fixed models, including the unmatched
L2 extension. Complete and incomplete columns have different denominators and
must not be added. Every $8/255$ positive count is zero. Three seeds do not
constitute a sample of all architectures, and the input-block intervals in the
main text do not quantify random-training uncertainty.

% Seed-count table inserted from the committed S1 final derived archive.
% S1 archive sha256 8dde94094d1fe97068ffa55dd32dd6204a8a9a911e089f292aa77c5f962adec5
\begin{table}[ht]
\centering\small
\caption{S1 numerical positives by fixed seed. C denotes complete verification (out of 200); I denotes incomplete CROWN-IBP (out of 1000 per model). No seed is omitted.}
\label{tab:s1seeds}
\begin{tabular}{llrrrr}
\toprule
Group & Seed & C ($2/255$) & I ($2/255$) & C ($8/255$) & I ($8/255$) \\
\midrule
L1E8 & 0 & 46 & 16 & 0 & 0 \\
L1E8 & 1 & 34 & 3 & 0 & 0 \\
L1E8 & 2 & 34 & 4 & 0 & 0 \\
L2E8 & 0 & 4 & 0 & 0 & 0 \\
L2E8 & 1 & 7 & 0 & 0 & 0 \\
L2E8 & 2 & 6 & 0 & 0 & 0 \\
Dense total & 0 & 26 & 0 & 0 & 0 \\
Dense total & 1 & 47 & 1 & 0 & 0 \\
Dense total & 2 & 42 & 1 & 0 & 0 \\
Dense active & 0 & 40 & 1 & 0 & 0 \\
Dense active & 1 & 35 & 1 & 0 & 0 \\
Dense active & 2 & 41 & 0 & 0 & 0 \\
\bottomrule
\end{tabular}
\end{table}

The committed S1 archive is
\begin{quote}\small
\path{docs/architecture_s1_result_20261031.json}.
\end{quote}
It contains 48 model/mode/radius rows, all eight paired comparisons, gained/lost input-seed
identities, and costs. Its audit hash binds the full original terminal records.
The eight old public-PGD errors are retained, accounting for eighty unaccepted
domains; they are neither rerun nor imputed as successful attacks. All 4800
complete and 24000 incomplete verification requests are settled without missing
costs. Shared-clean-correct strata explain all-request timing but do not alter
the registered primary denominator. S2 remains unexecuted after the stop.

The existing table script is
\begin{quote}\small
\path{scripts/render_current_review_tables.py}.
\end{quote}
Running it with \texttt{python -I -S -B} and \texttt{--check} reconstructs the
seven historical tables from their committed source summaries.
It does not load a model, execute an attack, or independently verify a numerical
bound. The current manuscript build separately checks the S1 transcription
against its final archive. Table reconstruction is not clean-environment
empirical reproduction; a human reviewer must distinguish those tasks.

\section{Source contracts and independent proof boundary}
\label{app:sourceboundary}
The historical output-layer ACT study reports 179 policy acceptances against
156 for matched monolithic and 141 for the older strong path, on 100 inputs
and three models. A later audit identifies input-containment gaps affecting
all 23 primary gains. Those counts remain numerical-policy outcomes; they are
not source-complete certificates. The requested and reconstructed input sets
are not nested. This does not establish that the models are unsafe, or predict
the results of a corrected execution. No sealed case is reopened here.

For a supplied LP $Ax\leq b$, $l\leq x\leq u$, objective $c^\top x+c_0$, and
any $y\leq0$, set $r=c-A^\top y$. Exact arithmetic can check the bound
\[
 c_0+b^\top y+\sum_i\min\{r_i l_i,r_i u_i\}.
\]
Indeed $y^\top Ax\geq y^\top b$, and the residual term is minimized over the
variable box. Multiplier signs, original coefficients, residual compensation,
property identity and obligation coverage must all be checked. This proves a
statement about the supplied LP. It does not establish input inclusion,
network-to-HZ conversion, guard lowering, omitted-history infeasibility or
native floating-point execution. A positive bound from an old matrix cannot
be attached to a new checked source construction.

Human technical review, a clean empirical reproduction, authorship and any
distribution and submission decisions remain external tasks. The manuscript and
record audits do not fill them or certify that the contribution meets a
particular venue's novelty standard. All unsuccessful gates, original errors,
and source limitations remain visible alongside local positive results.

\end{document}
